\chapter{Negative frequencies and connected ground correlations}
\label{chap:negative_frequencies}

This chapter proves the negative-frequency step of the exponential
clustering theorem of Hastings and Koma
\cite[Section 3, Lemma 3.1]{Hastings2006Spectral}. For a Hermitian
matrix with a spectral gap above a one-dimensional ground space, the
connected ground correlation of two observables is the ground expectation
of a commutator with the energy-lowering part of the first. A Gaussian
regularization of that part has a commutator bound uniform in the
regularizer whenever the Heisenberg evolution satisfies a Lieb--Robinson
bound, and its ground expectations approximate the sharp ones up to a
Gaussian error in the gap. The Lieb--Robinson bound enters as a
hypothesis.

\section{Centering and the energy-lowering part}
\label{sec:hk_energy_lowering}

\begin{lemma}[Centering an observable at an unnormalized vector]
    \label{lem:hk_observable_centering}
    \lean{ManyBody.centeredObservable,
        ManyBody.norm_inner_div_norm_sq_le,
        ManyBody.inner_centeredObservable_eq_zero,
        ManyBody.norm_centeredObservable_le,
        ManyBody.inner_centeredObservable_mul,
        ManyBody.centeredObservable_mem,
        ManyBody.centeredObservable_commutator}
    \leanok
    For $\psi\ne0$, write
    $\langle A\rangle_\psi=\langle\psi,A\psi\rangle/\|\psi\|^2$
    and $A_0=A-\langle A\rangle_\psi\Id$. Then
    $|\langle A\rangle_\psi|\le\|A\|$,
    $\langle\psi,A_0\psi\rangle=0$, and $\|A_0\|\le2\|A\|$.
    The centered observable gives the connected correlation:
    \[
        \frac{\langle\psi,A_0B\psi\rangle}{\|\psi\|^2}
        =\langle AB\rangle_\psi
        -\langle A\rangle_\psi\langle B\rangle_\psi.
    \]
    Subtracting this scalar identity preserves every linear space of
    observables containing $A$ and $\Id$, as well as the commutator with $B$.
\end{lemma}
\begin{proof}\leanok
    Use Cauchy--Schwarz, the operator-norm bound, and the triangle
    inequality. Expand the scalar identity and use that $\|\psi\|^2>0$.
\end{proof}

\begin{definition}[The energy-lowering part of an observable]
    \label{def:hk_negative_frequency_part}
    \lean{ManyBody.negativeFrequencyPart}
    \leanok
    Let $H$ be a finite Hermitian Hamiltonian with an orthonormal
    eigenbasis $e_i$ and energies $E_i$. For an observable $A$, define
    \[
        A^- = \sum_{E_i<E_j}
        \langle e_i,Ae_j\rangle\,|e_i\rangle\langle e_j|.
    \]
    This is the sharp negative-frequency part of the Heisenberg
    evolution $e^{itH}Ae^{-itH}$ in
    \cite[Section 3, equations (3.1)--(3.6)]{Hastings2006Spectral}.
\end{definition}

\begin{lemma}[Exact separation of the connected correlation]
    \label{lem:hk_negative_frequency_ground_correlation}
    \lean{ManyBody.negativeFrequencyPart_apply_ground_eq_zero,
        ManyBody.inner_ground_projected_eq_negativeFrequency_commutator,
        ManyBody.inner_ground_connected_eq_negativeFrequency_commutator}
    \leanok
    \uses{def:hk_negative_frequency_part,def:hamiltonians_gap}
    Suppose $E_0\le E_i$ for every $i$, let $\psi$ belong to the
    eigenspace at $E_0$, and let $P_0$ project onto that full eigenspace.
    Then $A^-\psi=0$ and
    \[
        \langle\psi,A(1-P_0)B\psi\rangle
        =\langle\psi,[A^-,B]\psi\rangle.
    \]
    If the ground eigenspace is $\mathbb C\psi$, this is
    \[
        \langle\psi,AB\psi\rangle-
        \frac{\langle\psi,A\psi\rangle\langle\psi,B\psi\rangle}
        {\|\psi\|^2}
        =\langle\psi,[A^-,B]\psi\rangle.
    \]
\end{lemma}
\begin{proof}\leanok
    An energy-lowering matrix element cannot act nontrivially on a
    ground vector. Expand an arbitrary vector orthogonal to the ground
    eigenspace in the energy basis. Between that vector and a ground
    bra, every nonzero matrix element of $A$ lowers energy, so it is
    retained by $A^-$. Apply this to $(1-P_0)B\psi$.
    The term $BA^-\psi$ vanishes. If the ground eigenspace is
    $\mathbb C\psi$, then
    \[
        P_0x=\frac{\langle\psi,x\rangle}{\|\psi\|^2}\,\psi,
        \qquad
        \langle\psi,AP_0B\psi\rangle
        =\frac{\langle\psi,A\psi\rangle\langle\psi,B\psi\rangle}{\|\psi\|^2},
    \]
    and substituting this into the first identity gives the second.
\end{proof}

The sharp energy-lowering part need not be local. The following lemmas
replace it by a Gaussian regularization, control that regularization by a
time integral of local commutators, and derive exponential clustering.

\section{The Gaussian--Cauchy transform}
\label{sec:hk_gaussian_cauchy}

\begin{lemma}[The Gaussian negative-frequency selector]
    \label{lem:hk_gaussian_negative_frequency_cutoff}
    \lean{GaussianFilter.gaussianNegativeFrequencyCutoff,
        GaussianFilter.gaussianNegativeFrequencyCutoff_nonneg,
        GaussianFilter.gaussianNegativeFrequencyCutoff_le_one,
        GaussianFilter.gaussianNegativeFrequencyCutoff_le_exp_of_gap,
        GaussianFilter.abs_one_sub_gaussianNegativeFrequencyCutoff_le_exp_of_gap}
    \leanok
    For $q,\Delta>0$, let $X$ be a centered Gaussian of variance
    $\Delta^2/q$ and put $F_{q,\Delta}(\omega)=\mathbb P(X\ge\omega)$.
    Then $0\le F_{q,\Delta}(\omega)\le1$ and
    \[
        \omega\ge\Delta\ \Longrightarrow\
        F_{q,\Delta}(\omega)\le e^{-q/2},
        \qquad
        \omega\le-\Delta\ \Longrightarrow\
        |1-F_{q,\Delta}(\omega)|\le e^{-q/2}.
    \]
    This is the Gaussian selector of
    \cite[Lemma 3.1, equation (3.7)]{Hastings2006Spectral}, with
    $\alpha=\Delta^2/(2q)$.
\end{lemma}
\begin{proof}\leanok
    The moment-generating function is
    $\mathbb E(e^{tX})=\exp(\Delta^2t^2/(2q))$.
    Chernoff's bound at $\Delta$ is $e^{-q/2}$. The positive-frequency
    statement follows by inclusion of the corresponding tail events.
    The complementary event for $\omega\le-\Delta$ is contained in
    $\{-X\ge\Delta\}$, and $-X$ has the same sub-Gaussian bound.
\end{proof}

The following regularized time-integral argument identifies this actual
Gaussian selector. The operator locality estimates are proved after it.

\begin{definition}[The regularized Gaussian--Cauchy kernel]
    \label{def:hk_gaussian_cauchy_kernel}
    \lean{GaussianFilter.gaussianCauchyKernel}
    \leanok
    For $\alpha,\varepsilon>0$, set
    \[
        K_{\alpha,\varepsilon}(t)=
        \frac{e^{-\alpha t^2}}{2\pi(\varepsilon-it)}.
    \]
    This is the regularized negative-frequency time kernel of
    \cite[Lemma 3.1, equation (3.8)]{Hastings2006Spectral}.
\end{definition}

\begin{lemma}[Laplace representation and integrability of the kernel]
    \label{lem:hk_gaussian_cauchy_integrability}
    \lean{GaussianFilter.cauchyDenominator_ne_zero,
        GaussianFilter.norm_inv_cauchyDenominator_le,
        GaussianFilter.inv_cauchyDenominator_eq_integral,
        GaussianFilter.norm_gaussianCauchyKernel_le,
        GaussianFilter.integrable_gaussianCauchyKernel_mul_exp,
        GaussianFilter.norm_integral_gaussianCauchyKernel_mul_exp_le}
    \leanok
    \uses{def:hk_gaussian_cauchy_kernel}
    For $\alpha,\varepsilon>0$ and $\omega,t\in\mathbb R$,
    \[
        \frac1{\varepsilon-it}
        =\int_0^\infty e^{(-\varepsilon+it)s}\,ds,
        \qquad
        |K_{\alpha,\varepsilon}(t)|
        \le\frac{e^{-\alpha t^2}}{2\pi\varepsilon}.
    \]
    The function $K_{\alpha,\varepsilon}(t)e^{i\omega t}$ is absolutely
    integrable and
    \[
        \left|\int_{\mathbb R}
        K_{\alpha,\varepsilon}(t)e^{i\omega t}\,dt\right|
        \le \frac1{2\pi\varepsilon}\sqrt{\frac\pi\alpha}.
    \]
\end{lemma}
\begin{proof}\leanok
    The complex exponential integral converges since its exponent has
    real part $-\varepsilon<0$. Its antiderivative gives the Laplace
    identity. The denominator has norm at least $\varepsilon$, so
    the kernel is dominated by the displayed Gaussian. The oscillatory
    exponential has norm one; Gaussian integrability and the Gaussian
    integral formula finish the proof.
\end{proof}

The bound retains its dependence on $\varepsilon$. It is not a uniform
regularizer-limit estimate. The next lemma interchanges the two integrals,
identifies the Gaussian selector, and passes to the regularizer limit.

\begin{lemma}[The scalar Gaussian--Cauchy limit]
    \label{lem:hk_gaussian_cauchy_limit}
    \lean{GaussianFilter.gaussianCauchyTransform,
        GaussianFilter.gaussianCauchyTransform_eq_integral,
        GaussianFilter.tendsto_gaussianCauchyTransform_zero_regularizer,
        GaussianFilter.halfLine_gaussian_integral_eq_measureReal,
        GaussianFilter.tendsto_normalized_gaussianCauchyTransform,
        GaussianFilter.tendsto_normalized_gaussianCauchyTransform_cutoff}
    \leanok
    \uses{lem:hk_gaussian_cauchy_integrability,lem:hk_gaussian_negative_frequency_cutoff}
    Let $\nu_h$ be the centered Gaussian probability measure of variance
    $h>0$, and let $c_h$ be the density at zero of $\nu_{1/h}$.
    For $\varepsilon>0$, define the transform
    \[
        C_{h,\varepsilon}(\omega)=\int_{\mathbb R}
        \frac{e^{i\omega t}}{\varepsilon-it}\,d\nu_h(t).
    \]
    Then
    \[
        C_{h,\varepsilon}(\omega)=
        \int_0^\infty e^{-\varepsilon s}
        e^{-h(\omega+s)^2/2}\,ds,
        \qquad
        c_hC_{h,\varepsilon}(\omega)
        \longrightarrow \nu_{1/h}([\omega,\infty))
    \]
    as $\varepsilon\downarrow0$.
    In particular, $h=q/\Delta^2$ gives the Gaussian selector
    $F_{q,\Delta}$ above and its explicit gap-error bounds.
\end{lemma}
\begin{proof}\leanok
    Insert the half-line Laplace representation of the Cauchy
    denominator. Under $\nu_h\otimes ds$ on $s>0$, the absolute value of
    the joint integrand is $e^{-\varepsilon s}$, which is integrable.
    Fubini and the exact Gaussian characteristic function give the first
    identity. For the limit, the translated Gaussian
    $e^{-h(\omega+s)^2/2}$ is an integrable majorant independent of
    $\varepsilon>0$. Apply dominated convergence.
    The density formula for $\nu_{1/h}$, translation invariance
    of Lebesgue integration, and equality of open and closed integral
    endpoints give the normalization to the probability tail.
\end{proof}

The scalar frequency selector and its integral representation supply
the regularized observable, ground-vector comparison, and locality
argument below.

\section{The regularized negative-frequency observable}
\label{sec:hk_regularized_observable}

\begin{definition}[The regularized negative-frequency observable]
    \label{def:hk_regularized_negative_frequency_operator}
    \lean{ManyBody.gaussianNegativeFrequencyIntegrand,
        ManyBody.gaussianNegativeFrequencyOperator}
    \leanok
    \uses{lem:hk_gaussian_cauchy_limit}
    For a Hermitian $H$, an observable $A$, Gaussian time variance $h$,
    and positive regularizer $\varepsilon$, put
    \[
        A^-_{h,\varepsilon}
        =c_h\int_{\mathbb R}
        \frac{e^{itH}Ae^{-itH}}{\varepsilon-it}\,d\nu_h(t).
    \]
\end{definition}

\begin{lemma}[regularized energy matrix elements]
    \label{lem:hk_regularized_frequency_matrix_elements}
    \lean{ManyBody.continuous_gaussianNegativeFrequencyIntegrand,
        ManyBody.norm_gaussianNegativeFrequencyIntegrand_le,
        ManyBody.integrable_gaussianNegativeFrequencyIntegrand,
        ManyBody.norm_gaussianNegativeFrequencyOperator_le,
        ManyBody.dotProduct_integral_gaussianNegativeFrequencyIntegrand_mulVec,
        ManyBody.dotProduct_gaussianNegativeFrequencyOperator_mulVec,
        ManyBody.inner_gaussianNegativeFrequencyOperator,
        ManyBody.tendsto_dotProduct_gaussianNegativeFrequencyOperator}
    \leanok
    \uses{def:hk_regularized_negative_frequency_operator,lem:hk_gaussian_cauchy_limit}
    The regularized observable is integrable and
    $\|A^-_{h,\varepsilon}\|\le c_h\varepsilon^{-1}\|A\|$.
    For energy eigenvectors $Hv=Ev$ and $Hw=E'w$,
    \[
        \langle v,A^-_{h,\varepsilon}w\rangle
        =c_hC_{h,\varepsilon}(E-E')\langle v,Aw\rangle.
    \]
    If $h>0$, the entries therefore converge as
    $\varepsilon\downarrow0$ to
    \[
        \nu_{1/h}([E-E',\infty))\langle v,Aw\rangle.
    \]
\end{lemma}
\begin{proof}\leanok
    Unitary conjugation preserves the norm, and the rational factor has
    norm at most $\varepsilon^{-1}$, so the finite Gaussian measure
    gives integrability. Between energy eigenvectors the actual
    Heisenberg evolution contributes $e^{it(E-E')}$. Passing the matrix
    element through the integral gives the scalar transform.
    Apply its proved Gaussian probability-tail limit.
\end{proof}

\begin{lemma}[Uniform locality after cancelling the Cauchy singularity]
    \label{lem:hk_regularized_frequency_uniform_locality}
    \lean{ManyBody.norm_commutator_gaussianNegativeFrequencyOperator_le}
    \leanok
    \uses{def:hk_regularized_negative_frequency_operator}
    Suppose $C,v\ge0$ and the dynamics satisfy
    $\|[A(t),B]\|\le C(e^{v|t|}-1)$ for all real $t$.
    For $T,\varepsilon>0$,
    \[
        \|[A^-_{h,\varepsilon},B]\|
        \le c_h\left(Cve^{vT}
        +\frac{2\|A\|\|B\|}{T}e^{-T^2/(2h)}\right).
    \]
    In particular the bound is independent of the positive regularizer.
\end{lemma}
\begin{proof}\leanok
    The rational factor is at most $|t|^{-1}$ for $t\ne0$.
    For $|t|\le T$, the inequality $e^x-1\le xe^x$ cancels this
    factor against the short-time increment. At $t=0$ the
    commutator vanishes. Outside the time window, the rational factor
    is at most $T^{-1}$ and the commutator norm at most
    $2\|A\|\|B\|$. Split the Gaussian integral and apply its
    Gaussian tail bound. Both bounds are independent of $\varepsilon$.
\end{proof}

\begin{lemma}[Gaussian ground-vector separation]
    \label{lem:hk_gaussian_ground_frequency_separation}
    \lean{ManyBody.exists_gaussianNegativeFrequency_ground_limits,
        ManyBody.exists_gaussianNegativeFrequency_ground_commutator_limit}
    \leanok
    \uses{lem:hk_regularized_frequency_matrix_elements,lem:hk_gaussian_negative_frequency_cutoff,
        def:hamiltonians_gap}
    Suppose the full ground eigenspace of $H$ is $\mathbb C\psi$,
    every other energy is at least $\Delta>0$ above it, and
    $\langle\psi,A\psi\rangle=0$. For $q>0$, put $h=q/\Delta^2$.
    The vectors $A^-_{h,\varepsilon}\psi$ and
    $(A^-_{h,\varepsilon})^\dagger\psi$ have limits $\xi,\eta$ with
    \[
        \|\xi\|\le e^{-q/2}\|A\psi\|,
        \qquad \|A^\dagger\psi-\eta\|
        \le e^{-q/2}\|A^\dagger\psi\|.
    \]
    The scalar commutator therefore converges to a number $z$ satisfying
    \[
        |\langle\psi,AB\psi\rangle-z|
        \le 2e^{-q/2}\|A\|\|B\|\|\psi\|^2.
    \]
    Here $\psi$ need not be a unit vector, and $B$ is arbitrary.
\end{lemma}
\begin{proof}\leanok
    Let $b_i$ be an orthonormal energy basis, $Hb_i=E_ib_i$, and write
    $F=F_{q,\Delta}$. The matrix-element limit gives
    \[
        \langle b_i,\xi\rangle=F(E_i-E_0)\langle b_i,A\psi\rangle,
        \qquad
        \langle b_i,\eta\rangle=F(E_0-E_i)\langle b_i,A^\dagger\psi\rangle.
    \]
    If $E_i=E_0$, then $b_i\in\mathbb C\psi$, and centering gives
    $\langle b_i,A\psi\rangle=\langle b_i,A^\dagger\psi\rangle=0$.
    Otherwise $E_i-E_0\ge\Delta$, so $F(E_i-E_0)\le e^{-q/2}$ and
    $|1-F(E_0-E_i)|\le e^{-q/2}$. By Parseval,
    \[
        \|\xi\|^2=\sum_i|\langle b_i,\xi\rangle|^2
        \le e^{-q}\sum_i|\langle b_i,A\psi\rangle|^2
        =e^{-q}\|A\psi\|^2,
    \]
    and in the same way
    $\|A^\dagger\psi-\eta\|^2\le e^{-q}\|A^\dagger\psi\|^2$.
    Finally,
    \[
        \langle\psi,AB\psi\rangle-
        \langle\psi,(A^-B-BA^-)\psi\rangle
        \longrightarrow
        \langle A^\dagger\psi-\eta,B\psi\rangle
        +\langle B^\dagger\psi,\xi\rangle,
    \]
    and Cauchy--Schwarz bounds each term by
    $e^{-q/2}\|A\|\|B\|\|\psi\|^2$.
\end{proof}


\begin{lemma}[Ordinary centered correlation from Gaussian frequency separation]
    \label{lem:hk_gaussian_ordinary_centered_correlation}
    \lean{ManyBody.norm_centered_ground_correlation_le_gaussian}
    \leanok
    \uses{lem:hk_gaussian_ground_frequency_separation,lem:hk_regularized_frequency_uniform_locality}
    Under the preceding ground-line and propagation hypotheses, for
    $q,T>0$ and $h=q/\Delta^2$,
    \[
        |\langle\psi,AB\psi\rangle|
        \le\left(2e^{-q/2}\|A\|\|B\|
        +c_h\left(Cve^{vT}
            +\frac{2\|A\|\|B\|}{T}e^{-T^2/(2h)}\right)\right)\|\psi\|^2.
    \]
\end{lemma}
\begin{proof}\leanok
    The regularized commutator is bounded uniformly in its positive
    regularizer. Pass this bound to its scalar limit and add the proved
    Gaussian ground-vector error.
\end{proof}

\begin{lemma}[An exponential rate for the Gaussian frequency bound]
    \label{lem:hk_gaussian_frequency_clustering_rate}
    \lean{GaussianFilter.kernel_inv_zero_eq_sqrt,
        GaussianFilter.exists_gaussianFrequencyClusteringRate}
    \leanok
    If $\Delta,v>0$, there are $C,\kappa>0$ such that for $m\ge0$ and
    $\ell\ge1$,
    \[
        2e^{-\ell\Delta/(4v)}+
        \sqrt{\frac{\ell}{4\pi v\Delta}}
        \left(2mv e^{-\ell/2}
        +\frac{4v}{\ell}e^{-\ell\Delta/(4v)}\right)
        \le C(m+1)e^{-\kappa\ell}.
    \]
\end{lemma}
\begin{proof}\leanok
    The density normalization is $c_h=\sqrt{h/(2\pi)}$.
    Take $\kappa=\min(1/4,\Delta/(8v))$ and absorb the square-root
    factors using $\sqrt\ell\le\ell$ and
    $\ell e^{-\ell/2}\le4e^{-\ell/4}$.
\end{proof}
